\section{From one apex to uniform inequalities}
\label{sec:wall}

The fixed-apex estimate now supplies the global inequality.  The
mechanism is numerical: a differential inequality preserves a strict
violation as the volume parameter moves toward the apex.  If an object
becomes strictly semistable, one continues with a violating stable
factor.  A discrete discriminant makes this process finite.
Transport along a zero-slope hyperbola and induction on the
discriminant appear in \cite[Section~5]{BMS2016} and, for a corrected
inequality, \cite[Lemma~2.7]{BMSZ2017}.  Here the terminal estimate is
at a fixed positive volume $A$.  The comparison functions below turn
that estimate into an all-volume correction and an uncorrected tail.

\subsection{A transport principle}

At slope zero define the defect and the nonnegative comparison quantity
\begin{equation}\label{eq:wall-defect}
 D_a(E)=z_b(E)-\frac{a^2d_b(E)}6,\qquad
 S_a(E)=\frac{\Delta(E)}{d_b(E)}
       =d_b(E)-\frac{a^2r(E)^2}{d_b(E)}.
\end{equation}
They are defined whenever $d_b(E)>0$ and $w_b(E)=a^2r(E)/2$.
In particular $0\le S_a\le d_b$ for a semistable object.

\begin{proposition}[Transport from a fixed apex]\label{prop:apex-transport}
Fix a smooth projective complex threefold $X$, an integral ample class
$H$, and $B_0\in N^1(X)_{\QQ}$.  Suppose that there are $A>0$ and
$C\ge0$ such that every finite-slope tilt-semistable object of slope
zero at $(A,b)$, for every $b\in\RR$, satisfies
\begin{equation}\label{eq:wall-apex-input}
 D_A(E)\le C S_A(E).
\end{equation}
Then, with
\begin{equation}\label{eq:wall-constants}
 R_0=A+\frac{3C}{A},\qquad
 k_0=\max\left\{C+\frac{A^2}{6},\frac{CR_0}{A}\right\},
\end{equation}
every finite-slope tilt-semistable object of slope zero at $(a,b)$
satisfies $D_a(E)\le k_0d_b(E)$ for all $a>0,b\in\RR$, and
$D_a(E)\le0$ whenever $a>R_0$.
\end{proposition}

\begin{proof}
We first compute along the slope-zero curve of a fixed class.  Its
positive-denominator branch is
\begin{equation}\label{eq:zero-branch}
 d(a)=\sqrt{\Delta+a^2r^2},\qquad
 b(a)=\frac{d_{B_0}-d(a)}r\quad(r\ne0).
\end{equation}
For rank zero, $d$ and $b$ remain constant.  If the branch starts at
a semistable object, then $\Delta\ge0$, and its denominator stays
positive at every positive $a$.  Indeed $d(a)\ge a|r|$ when
$r\ne0$, whereas $d$ is a positive constant when $r=0$.
The identities $dz_b/db=-w_b$ and $w_b=a^2r/2$ give
\begin{equation}\label{eq:wall-derivatives}
 D_a'=-\frac a3S_a,\qquad
 S_a'=-\frac{ar^2}{d(a)^2}S_a.
\end{equation}
For example, when $r\ne0$ one has $d'=ar^2/d$,
$b'=-ar/d$, and $z'=a^3r^2/(2d)$; these give both equalities
directly.  The rank-zero case follows from the definitions.
For a differentiable nonnegative function $k$, therefore,
\begin{equation}\label{eq:wall-comparison-derivative}
 (D_a-k(a)S_a)'=
 S_a\left(-\frac a3-k'(a)+k(a)\frac{ar^2}{d(a)^2}\right).
\end{equation}

There are three comparisons.  For the corrected bound below $A$,
take $k(a)=C+(A^2-a^2)/6$ and move upward toward
$A$.  The right side of \eqref{eq:wall-comparison-derivative} is
$k(a)ar^2S_a/d(a)^2\ge0$, so a strict violation persists upward.
To rule out positive defect at a starting value $R>R_0$, take
$k(a)=a(R-a)/3$ on $[A,R]$ and move downward.  Since
$ar^2/d(a)^2\le1/a$,
\[
 -\frac a3-k'(a)+k(a)\frac{ar^2}{d(a)^2}
 \le-\frac a3-k'(a)+\frac{k(a)}a=0.
\]
A strict violation again persists in the direction of movement.
Here $k(R)=0$ and $k(A)=A(R-A)/3>C$.
For the corrected bound in the remaining interval, on any
interval $[A,R]$ use $k(a)=Ca/A$ and move
downward.  The same upper bound for the last term gives a derivative
at most $-aS_a/3\le0$, and $k(A)=C$.

Each comparison preserves a strict violation while the class remains
stable.  We next show how to continue through a change of stability,
with a strict decrease of a discrete discriminant at each replacement.
For an equal-slope-zero filtration with stable factors $E_i$,
all $d_i$ are positive and $w_i=a^2r_i/2$.  Additivity and
weighted Cauchy--Schwarz give
\begin{equation}\label{eq:wall-additivity}
 D_a(E)=\sum_iD_a(E_i),\qquad
 S_a(E)=\sum_i d_i-a^2\frac{(\sum_i r_i)^2}{\sum_i d_i}
        \ge\sum_iS_a(E_i).
\end{equation}
Thus a strict inequality $D_a(E)>k(a)S_a(E)$, with $k(a)\ge0$,
passes to at least one stable factor.  Such a factorization is finite
by Lemma~\ref{lem:finite-tilt-factors}.

Start with a violating stable factor and follow its branch
\eqref{eq:zero-branch} toward $A$ as long as it remains stable.
Stability is open.  At an endpoint in the positive $a$-interval,
continuity of extremal phases gives semistability, since the formal
denominator remains positive.  The strict violation persists at the
endpoint: along the movement $D-kS$ is at least its initial positive
value.  If the endpoint is $A$, this contradicts
\eqref{eq:wall-apex-input}.  Otherwise the endpoint is nonstable,
since stability there would extend the interval, and we split again
and select a violating stable factor.

At every such genuine wall, the discriminant of each factor is
strictly smaller than that of its parent.  Indeed,
\begin{equation}\label{eq:wall-discriminant-split}
 \Delta(E)=\sum_i\Delta(E_i)+
    2\sum_{i<j}(d_i d_j-a^2r_i r_j).
\end{equation}
Every term is nonnegative because $d_i\ge a|r_i|$.
If all cross terms vanished, all $r_i$ would be nonzero with the
same sign and $d_i=a|r_i|$; hence all factors would have
$\Delta(E_i)=0$ and proportional triples $(r_i,d_i,w_i)$.
Twisting preserves this proportionality.  Openness keeps the
finitely many stable factors stable near the wall, so their fixed
extension triangles make $E$ an extension of objects with equal
finite slope immediately before the wall as well.  The filtration
has at least two factors, so this contradicts the preceding stability
of $E$.  Some cross term in \eqref{eq:wall-discriminant-split} is
therefore positive, and every factor has smaller discriminant.

Because $B_0$ is rational, formula
\eqref{eq:projected-discriminant} puts all these discriminants in
$N^{-1}\ZZ$ for one positive integer $N$ depending only on
$X,H,B_0$.  They are nonnegative on the semistable factors.
Strict descent consequently permits only finitely many wall splits.
This also excludes an accumulating sequence of walls with changing
factors: every change consumes at least $1/N$ of discriminant, and
every final stable interval either reaches $A$ or has an interior
semistable endpoint at which another such split is necessary.
There is no denominator-zero endpoint inside the interval, by
\eqref{eq:zero-branch}.  The proposed violation must therefore reach
$A$, giving a contradiction in each of the three comparisons.

The first comparison proves
$D_a\le(C+(A^2-a^2)/6)S_a$ for $0<a\le A$.
The second proves $D_a\le0$ for $a>R_0$.
The third proves $D_a\le(Ca/A)S_a$ for $A\le a\le R_0$.
Finally use $S_a\le d_b$ and \eqref{eq:wall-constants}.
\end{proof}

\subsection{The corrected inequality and the uncorrected tail}

\begin{proof}[Proof of Theorem~\ref{thm:uniform-inequality}]
First use the very-ample and transverse-twist reductions of
Section~\ref{sec:tilt}.  Proposition~\ref{prop:apex} supplies
$A>0$ and $C\ge0$, chosen independently of $E$ and $b$, such that
\[
 D_A(E)\le C\bigl(d_B(E)-A|r(E)|\bigr).
\]
At slope zero put $e=d_B-A|r|$.  Since $d_B\ge A|r|$,
\[
 S_A=\frac{(d_B-A|r|)(d_B+A|r|)}{d_B},\qquad
 e\le S_A\le2e.
\]
The hypothesis of Proposition~\ref{prop:apex-transport} follows.
Multiplying its conclusion by $h$ gives
\[
 \ch_3^B(E)\le\frac{a^2}{6}H^2\ch_1^B(E)
                      +k_0H^2\ch_1^B(E),
\]
and gives the same inequality without the correction for $a>R_0$.
Undo the polarization rescaling as described in Section~\ref{sec:tilt}.
This gives constants $k\ge0$ and $R_*>0$ for the original
$X,H,B_0$, uniform in both $E$ and $b$.

Enlarge $k$ to a rational nonnegative number.  The curve class
$\Gamma=kH^2\in N_1(X)_{\QQ}$ then satisfies
$\Gamma\cdot H=kh\ge0$, and the correction is exactly
$\Gamma\cdot\ch_1^B(E)$.  This proves the corrected statement
with its required rational class, as well as the independent
uncorrected large-volume threshold.  No canonical-bundle hypothesis
has entered either argument.
\end{proof}

\subsection{The quadratic inequality at every slope}

We finish by expressing the uncorrected tail in the alternative
coordinates used for generalized Bogomolov--Gieseker inequalities.
Here $B_0=0$ and
\[
 v(E)=(v_0,v_1,v_2,v_3)
 =(H^3\ch_0(E),H^2\ch_1(E),H\ch_2(E),\ch_3(E)).
\]
For $\alpha>\beta^2/2$, the first tilt is
$\mathcal B_{H,\beta H}$ and its slope is
\begin{equation}\label{eq:alternative-slope}
 \widehat\nu_{\alpha,\beta}(E)=
 \frac{v_2(E)-\beta v_1(E)+(\beta^2-\alpha)v_0(E)}
      {v_1(E)-\beta v_0(E)},
\end{equation}
with value $+\infty$ at zero denominator.  Define
\begin{equation}\label{eq:quadratic-form}
 Q_{\alpha,\beta}(v)=
 \alpha(v_1^2-2v_0v_2)+\beta(3v_0v_3-v_1v_2)
                      +2v_2^2-3v_1v_3.
\end{equation}
Thus no correction of any character is being incorporated into $v$.

\begin{proof}[Proof of Corollary~\ref{cor:quadratic-tail}]
Take the threshold of Theorem~\ref{thm:uniform-inequality} with
$B_0=0$.  Set $b=\beta$ and $a=\sqrt{2\alpha-b^2}$.
The numerator in \eqref{eq:alternative-slope} is
$h(w_b-a^2r/2)$ and its denominator is $hd_b$, so its slope is
exactly $\nu_{a,b}$.  Expanding \eqref{eq:quadratic-form} with
\eqref{eq:twist-two}--\eqref{eq:twist-three} gives
\begin{equation}\label{eq:quadratic-twisted}
 \frac{Q_{\alpha,b}(v(E))}{h^2}
    =\frac{a^2\Delta(E)}2+2w_b(E)^2-3d_b(E)z_b(E).
\end{equation}
If $d_b(E)=0$, this is nonnegative by
\eqref{eq:tilt-discriminants}; it includes zero-dimensional objects.

Otherwise put $u=\nu_{a,b}(E)$, $c=b+u$, and
$R=\sqrt{a^2+u^2}$.  Lemma~\ref{lem:tilt-semicircle} makes $E$
semistable of slope zero at $(R,c)$, with $d_c>0$.
Using $w_b=ud_b+a^2r/2$,
\[
 d_c=d_b-ur,\qquad
 z_c=z_b-uw_b+\tfrac12u^2d_b-\tfrac16u^3r
\]
in \eqref{eq:quadratic-twisted} gives the exact identity
\begin{equation}\label{eq:quadratic-zero-slope-identity}
 \frac{Q_{\alpha,b}(v(E))}{h^2}
       =-3d_b(E)\left(z_c(E)-\frac{R^2d_c(E)}6\right).
\end{equation}
Since $\alpha>f(b)$ implies $R\ge a>R_*$, the expression in
parentheses is nonpositive by the uncorrected tail.  This proves the
claim.  The passage uses the entire semistability arc and the
inequality $R\ge a$; an all-parameter equivalence alone would not
justify a conclusion on this truncated region.
\end{proof}
